\songtitle{Trois fois prisonnier}

\tag*{2026}\tag{hip-hop}\tag{french-lyrics}\tag{african-empire}\tag{monukae}\tag{abulon}\tag{thrice-a-prisoner}
\begin{notes}
French-language adaptation by Claude of \emph{Thrice a Prisoner}, from the \emph{Monukae} saga within the \emph{African Empire} alternate-history universe by Carlos Thompson. Retrieved from rataflechera's Suno page (V6, 13 September 2026).
\end{notes}
\begin{style}
Hip-hop trap épico-cinématique en marche half-time, baritons français parlés sur basses distordues et cordes ostinato, tam-taïko et djembé fusionnés avec motifs kora et flûte andine, chœurs graves, cuivres larges, ruptures tribales, tape saturation, plate reverb et ascensions orchestrales sombres
\end{style}
\begin{lyrics}
\annotation{Couplet 1}\\
La première fois, il fut prisonnier,\\
Abulon, gamin, à peine formé,\\
Un fort dans le désert, au pied des montagnes,\\
Dans la nuit, l'assaut soudain les gagne.

Les soldats de l'Inca trient les captifs,\\
Moitié la gorge tranchée, plus rien de vif,\\
Moitié le fond des mines, un sort massif,\\
Mais Abulon portait un éclat furtif.

Abulon fut envoyé au palais,\\
Armure de fer, drapé de soie parfait,\\
Trophée de guerre pour l'Inca exhibé,\\
Un levier prêt à être actionné.

\annotation{Refrain}\\
Trois fois prisonnier\\
Trois fois vainqueur\\
L'adversité comme plan\\
Naissance d'une légende\\
Trois fois prisonnier\\
Aucune geôle ne l'a dompté\\
L'homme le plus puissant sur terre\\
Ainsi conte la légende

\annotation{Couplet 2}\\
L'Inca voyait en lui un simple levier,\\
Mais Abulon a su le déjouer,\\
Soldat stratège, diplomate rusé,\\
Quatre ans plus tard, tout a basculé :\\
L'Inca devint son beau-père fidèle,\\
Un traité lie commerce et citadelles.

Seconde fois, il fut prisonnier,\\
À peine revenu, à peine posé,\\
Des usurpateurs ont pris la terre convoitée,\\
Monukaé assiégée, son père achevé.

Avec les lances incas, les épées de ses frères,\\
Abulon déchaîne sa fureur, sa colère,\\
Il libère sa mère, mais tombe en la bataille dernière,\\
Traîné devant l'empereur — accusé, traître.

\annotation{Refrain}\\
Trois fois prisonnier\\
Trois fois vainqueur\\
L'adversité comme plan\\
Naissance d'une légende\\
Trois fois prisonnier\\
Aucune geôle ne l'a dompté\\
L'homme le plus puissant sur terre\\
Ainsi conte la légende

\annotation{Couplet 3}\\
L'empereur voit ce que voyait l'Inca avant,\\
Un homme d'épée, du cœur du peuple, confiant,\\
Qui maîtrise et le pouvoir et la négociation,\\
Un allié précieux, un atout du moment.

Les terres d'Abulon rendues, à condition\\
De servir l'empereur, ses hommes en soumission,\\
Au service du Mali, loyal sans hésitation,\\
Il devient bientôt son plus proche conseiller,\\
Général par mérite, magnat du commerce prospère.

Mais l'empereur meurt, son fils est indigne,\\
Les nobles se réunissent pour élire un signe,\\
Abulon s'impose, choix évident dès le premier tour,\\
Mais menace de partir si on le couronne un jour.

\annotation{Refrain}\\
Trois fois prisonnier\\
Trois fois vainqueur\\
L'adversité comme plan\\
Naissance d'une légende\\
Trois fois prisonnier\\
Aucune geôle ne l'a dompté\\
L'homme le plus puissant sur terre\\
Ainsi conte la légende

\annotation{Couplet 4}\\
Général, noble, maître d'une flotte fière,\\
Qui vend le fer à l'Inca contre or et argent, sans manière,\\
Aujourd'hui régent, il veille sur le jeune souverain,\\
Et l'héritier désigné, gardien de son destin.

Le fils d'Abulon, Mombuto, pris en mer Cantabrique,\\
Emmené en Irlande contre rançon, tactique,\\
De Mac Ghildt — Abulon prépare une flotte de trésor,\\
Le double de la rançon, pour sceller l'accord.

Mac Ghildt crut la prise trop facile, trop belle,\\
Capture la flotte — croit tenir la ficelle :\\
Si Abulon cède si vite, il peut payer plus fort,\\
Pour son fils et son émissaire — sans savoir encore\\
Que l'émissaire, c'était Abulon lui-même, en personne.

\annotation{Refrain}\\
Trois fois prisonnier\\
Trois fois vainqueur\\
L'adversité comme plan\\
Naissance d'une légende\\
Trois fois prisonnier\\
Aucune geôle ne l'a dompté\\
L'homme le plus puissant sur terre\\
Ainsi conte la légende

\annotation{Pont}\\
Derrière l'horizon, cachée, la flotte au trésor,\\
Deux mille trois cents hommes attendaient l'essor,\\
Cork fut puni, mais Mac Ghildt s'échappa encore,\\
La ville rasée, et dans les cendres qui dorment,\\
L'or de la rançon — le double promis —\\
Fut laissé, en guise d'avis.

\annotation{Refrain}\\
Trois fois prisonnier\\
Trois fois vainqueur\\
L'adversité comme plan\\
Naissance d'une légende\\
Trois fois prisonnier\\
Aucune geôle ne l'a dompté\\
L'homme le plus puissant sur terre\\
Ainsi conte la légende

\songpart{Outro}
L'Europe comprit, au sud du Sahara,\\
Qu'un pouvoir nouveau désormais s'imposera,\\
Mac Ghildt fut enfin capturé par les Francs,\\
Emmené à Monukaé, gorge tranchée par son enfant.
\end{lyrics}

